% Part: sets-functions-relations
% Chapter: sets
% Section: important-sets

\documentclass[../../../include/open-logic-section]{subfiles}

\begin{document}

\olfileid{sfr}{set}{imp}
\olsection{Een paar belangrijke verzamelingen}

\begin{ex}
We werken meestal met verzamelingen waarvan de !!{element}en wiskundige
objecten zijn. Vier van die verzamelingen zijn zo belangrijk dat ze een eigen
naam hebben:
\begin{multline*}
    \Nat = \{0, 1, 2, 3, \ldots\} \\
    \shoveright{\text{de verzameling natuurlijke getallen}}\\
    \shoveleft{\Int = \{\ldots, -2, -1, 0, 1, 2, \ldots\}} \\
    \shoveright{\text{de verzameling gehele getallen}}\\
    \shoveleft{\Rat = \Setabs{\nicefrac{m}{n}}{m, n \in \Int\text{ en }n \neq 0}}\\
    \shoveright{\text{de verzameling rationale getallen}}\\
    \shoveleft{\Real = (-\infty, \infty)}\\
    \text{de verzameling reële getallen (het continuüm)}
\end{multline*}
Deze verzamelingen zijn allemaal \emph{oneindig}: elke verzameling bevat
oneindig veel !!{element}en.

Bij elke volgende verzameling komen er \emph{meer} getallen bij. Daarom is
$\Nat \subseteq \Int \subseteq \Rat \subseteq \Real$: elk natuurlijk
getal is ook een geheel getal, elk geheel getal is ook rationaal en elk
rationaal getal is ook reëel. Bovendien is $\Nat \subsetneq \Int
\subsetneq \Rat$. Het getal $-1$ is namelijk wel geheel maar niet natuurlijk,
en $\nicefrac{1}{2}$ is wel rationaal maar niet geheel. Dat $\Rat
\subsetneq \Real$ is minder vanzelfsprekend. Er zijn dus reële getallen die
niet rationaal zijn\oliflabeldef{sfr:arith:real:realline}{; daar komen we op
terug in \olref[arith][real]{realline}}{}.

We gebruiken soms ook de positieve gehele getallen $\PosInt = \{1,
2, 3, \dots\}$ en de verzameling met alleen de eerste twee natuurlijke
getallen, $\Bin = \{0, 1\}$.
\end{ex}

\begin{tagblock}{compsci}
\begin{ex}[Tekenreeksen]
Nog een voorbeeld is $A^{*}$, de verzameling van alle \emph{eindige
tekenreeksen} die je met een alfabet $A$ kunt maken. Elke eindige rij
elementen van~$A$ is zo'n tekenreeks over~$A$. Ook de
\emph{lege tekenreeks $\Lambda$},
waarin geen enkel teken staat, hoort bij de tekenreeksen over~$A$, voor elk
alfabet~$A$. Bijvoorbeeld:
\begin{multline*}
\Bin^*
=\{\Lambda,0,1,00,01,10,11,\\
000,001,010,011,100,101,110,111,0000,\ldots\}.
\end{multline*}
Als $x=x_{1}\ldots x_{n}\in A^{*}$ uit $n$ ``letters'' van $A$ bestaat,
dan is de \emph{lengte} van die tekenreeks~$n$. We schrijven dan
$\len{x}=n$.
\end{ex}
\end{tagblock}

\begin{ex}[Oneindige rijen]
Bij elke verzameling $A$ kunnen we ook kijken naar~$A^\omega$, de verzameling
van alle oneindige rijen !!{element}en van~$A$. Zo'n rij
$a_1a_2a_3a_4\dots$ is een lijst die aan één kant eindeloos doorloopt.
Elk object in de lijst is !!a{element} van~$A$.
\end{ex}
\end{document}
